% Chapter 11: The World of Rotation
\chapter{The World of Rotation}

\step{A Counterintuitive Opening}

\begin{mainline}
Two eggs sit on the breakfast table, one raw and one hard-boiled. They are equally large, heavy, and white; with your eyes closed, you cannot feel any difference. Now do something a three-year-old can do: lay each on the table and give it a spin with your fingers, like a top.

The difference is immediate. The boiled egg spins fast and steadily, a lively little top whirling for more than ten seconds. The raw egg is sluggish: after barely two or three turns, it wobbles to a halt, as though something inside the shell were holding it back.

That is not the strangest part. Spin the boiled egg again, then gently hold it with a finger for about one second, until it stops completely. Let go. You have apparently confiscated its rotation, yet the instant you release it, it starts spinning again! Who rewound the spring inside? An egg contains no engine or spring, only white and yolk.

Do not rush to the answer. There are many similar puzzles. A figure skater pulls in her outstretched arms and spins several times faster, though nobody pushes her. A bicycle resists falling when its wheels spin, but wobbles over when they stop. Earth has been rotating for over four billion years; nobody ever winds it up, yet it has never missed a day. Behind these phenomena is a set of laws as fundamental as those of pushing, but often overlooked: the laws of rotation. Our task is to uncover them in eggs, ice skates, and bicycle wheels.
\end{mainline}

\step{Make a Guess First}

As usual, place your bets before seeing the cards. Write down predictions and reasons for these three questions. Wrong answers cost no points: intuition is precisely what this chapter aims to repair.

First: the boiled egg ``revives'' after you stop and release it. Compared with before you stopped it, will its new spin be (a) equally fast; (b) noticeably slower; or (c) faster?

Second: a figure skater pulls in her arms and triples her spin rate. Where does the extra speed come from? (a) The ice gives her a push; (b) her muscles do work and speed her up; or (c) something is conserved and redistributed.

Third: push a closed door with the same force in three places: at the handle, farthest from the hinge; at the middle; and right beside the hinge. Which push opens it most easily? What happens if you push the handle along the plane of the door, straight toward the hinge?

Write down your answers. We will check them one by one at the end.

\step{Hands-on Experiment}

\begin{experiment}[An Egg Spin Test and a Twisting Ruler]
\textbf{Materials}: one raw and one boiled egg (label them so you do not crack the wrong one), a table; a long ruler or light wooden stick, half a meter to a meter long; two equal lumps of modeling clay or equally heavy balls of wet tissue; and tape.

\textbf{Step 1: Tell the eggs apart.} Lay each egg on its side and spin it with roughly the same flick, three times apiece. Record which spins longer and more steadily. Try the ``revival'' test on the longer-spinning egg: hold it for one second, release it, and see whether it spins again. Compare its speed before and after revival.

\textbf{Step 2: Twist the bare ruler.} Hold the ruler horizontally with both hands at its center and rapidly twist it back and forth, a little left and a little right, as though wringing a towel. How hard is it to change its rotational motion?

\textbf{Step 3: Add dumbbells and twist again.} Tape the clay lumps to the ruler's ends, then hold the middle and twist it again. With the same hands and the same motion, the effort changes immediately: your arms struggle to get it moving. Finally, move the clay right next to your hands, still on the ruler and without changing the total weight. Twist again: it becomes easy.

\textbf{Record}: the ruler's total weight was the same in all three trials. Why did starting, stopping, and reversing its rotation feel so different? Clay farther from your hands makes twisting harder: what role does that distance play? Write this down alongside your egg observations. We will explain each result later. What you have just felt is something physicists call ``moment of inertia.''
\end{experiment}

Here is another free experiment to try at any door. Open it by pushing at the handle with one finger. Then push beside the hinge with one finger and the same effort. The second time, you will feel like a baby who has not learned how doors work. Remember that feeling: the same force produces different effects. It is the starting point for this entire chapter.

\step{The Old Model Runs into Trouble}

Review our old tools. The models developed in Chapters 6--10 have been highly successful. Describe an object by its position and velocity; predict its future by calculating its net force, which determines how fast its momentum changes. Momentum is conserved through collisions, and energy through processes. This language handles carts, billiard balls, and rockets.

But rotation brings one difficulty after another.

First, force alone cannot describe rotation. In the door experiment, three pushes have exactly the same magnitude and direction, but very different effects. The old equation \(F=ma\) contains no information about where the force acts: it assumes that you are pushing a point. A door is not a point, but an extended object distributed around its hinge. To predict whether it rotates and how fast, knowing the force is insufficient; we must also know how far from the axis it acts. Half the information is missing.

Second, rotation has its own version of ``What keeps it going?'' In Chapter 6, a curling stone kept sliding after release, forcing us to recognize that uniform straight-line motion needs no support. Now ask the same question about rotation: Earth has spun for over four billion years without an external push. Why does it never tire? Our old intuition says spinning things must eventually slow to a stop. But if a state of motion needs no explanation and only a \emph{change} does, rotation should follow the same rule: with nothing to change it, it keeps spinning as before. Our old model has no concept in which to put that statement.

Third, the egg revives. You hold the boiled egg still for a second, but its rotation returns when you let go. Chapter 10 taught us that momentum neither appears nor disappears from nothing; it is transferred. The egg's revival suggests that its rotation never truly vanished during that second. It hid somewhere your finger could not stop: inside the shell. But what does ``hidden rotation'' mean? Our old model has an account for straight-line motion, momentum, but none for rotation. To understand the egg, we need a new ledger.

All three difficulties point to the same diagnosis: our language for describing the world is missing an entire vocabulary of rotation. What we lack, we invent. That is how physics has worked for centuries.

\step{Inventing New Concepts}

\begin{mainline}
\textbf{The first group: describing rotation with angular position, angular velocity, and angular acceleration.} For straight-line motion we give position, velocity, and acceleration. Rotation follows the same recipe. Every point on an object rotating about a fixed axis stays the same distance from that axis; what changes is how far around it has moved. We therefore assign a marked line on a turntable an \emph{angular position}: how far it has turned from its starting direction. We can count turns, but a more useful unit is the \emph{radian}. If the marked point is a distance \(r\) from the axis, divide the arc length it travels by \(r\) to obtain the angle in radians. The benefit is immediate: arc length equals radius times angle in radians. A point twice as far from the axis travels twice as far through the same angle; one number describes every point along the radius. A full turn covers the entire circumference and corresponds to \(2\pi\) radians, about \(6.28\).

The rate at which angular position changes is \emph{angular velocity}: how many radians turn past each second. The everyday measure ``revolutions per minute'' is a close relative. The rate at which angular velocity itself changes is \emph{angular acceleration}. These three quantities are the rotational relatives of position, velocity, and acceleration. Their family roles match: where are we now, how fast is that changing, and how fast is the rate itself changing?

A wall clock is a ready-made display of angular velocities. Its hour, minute, and second hands share an axis but turn at different rates: the second hand once per minute, the minute hand once per hour, and the hour hand once per half-day, in the ratio \(60:1:\tfrac1{12}\). More interestingly, points near the hub and at the tip of the same second hand have identical \emph{angular velocities}, each completing a turn per minute, but their \emph{linear velocities} differ by tens of times. The tip travels an arc of nearly ten centimeters each minute while the hub barely moves. Remember: rotating points share angular velocity but each has its own linear velocity. Many later misconceptions come from confusing the two.

\textbf{The second concept: torque. Force alone does not determine a force's rotational effect.} Our experience with doors and screws boils down to a product of two factors: the force's magnitude and the perpendicular distance from the axis to its line of action, called the \emph{lever arm}. We call this product \emph{torque}. Large torque changes rotation rapidly; zero torque leaves it unchanged. Think of the leverage you gain when prying something up: force is the raw material, and the lever arm is the multiplier. What is torque not? It is \emph{not} a new force. Nothing mysterious lurks beside the hinge. Torque is the combined effect of force and position, just as work in Chapter 9 combined force and distance. That explains why handles sit farthest from hinges, wrenches have long handles, and a steel pipe slipped over a wrench frees a stubborn nut. Engineers do not waste lever arms.

Look around: torque is built into designs everywhere. A thick screwdriver handle places your grip farther from the axis, making the same screwdriver easier to turn. Car steering wheels are much larger than game controllers' wheels, and bus wheels are larger still, for the same reason. Slip a spoon handle under a can's ring-pull and pry: the can's rim becomes the pivot. Even when wringing a towel, you instinctively spread your hands toward its ends instead of crowding them at the middle. People exploited this law long ago: a small sliding weight on a steelyard can weigh tens of kilograms merely by changing its position. Its scale marks lever arms, not weights.
\end{mainline}

\begin{center}
\begin{tikzpicture}[scale=1.0]
  % Top view: door extending right from its hinge
  \fill[pattern=north east lines] (-0.15,-0.9) rectangle (0.15,0.9);
  \node[left] at (-0.2,0) {Hinge};
  \draw[very thick,gray!60!black] (0,0) -- (5.6,0);
  \node[above] at (4.6,0.02) {Door, top view};
  % Force 1: push the handle perpendicular to the door
  \draw[-{Stealth[length=3mm]},very thick,red!70!black] (5.2,0) -- (5.2,1.6) node[above] {Force \(\to\) longest arm: best effect};
  \draw[fill] (5.2,0) circle (0.05);
  % Force 2: push at the middle
  \draw[-{Stealth[length=3mm]},very thick,orange!80!black] (2.8,0) -- (2.8,1.6) node[above,yshift=16pt] {Same force, half the arm: half the effect};
  \draw[fill] (2.8,0) circle (0.05);
  % Force 3: push beside the axis
  \draw[-{Stealth[length=3mm]},very thick,blue!70!black] (0.9,0) -- (0.9,-1.3) node[below] {Near-zero arm: cannot open it};
  \draw[fill] (0.9,0) circle (0.05);
\end{tikzpicture}
\end{center}

\begin{mainline}
\textbf{The third concept: moment of inertia. Resistance to a change in rotation depends on both mass and where it sits.} In the twisting experiment, the ruler and clay kept the same total weight; what changed was their distance from the axis. The farther out the weight sat, the harder it was to twist back and forth. Resistance to a change in rotation therefore depends not only on how much mass there is, but on where it is distributed. We name this resistance \emph{moment of inertia}. It is the rotational counterpart of mass, our measure in Chapter 6 of resistance to velocity changes, with an extra trait: \emph{mass farther from the axis has a larger say}, amplified by the square of the distance. Move the same clay twice as far out and its resistance quadruples. A skater with outstretched arms thus has a large moment of inertia and spins slowly; pulling in her arms sharply reduces it. Likewise, a small solid mooncake tin and a large hollow iron ring of equal weight behave very differently when spun.

Toy designers exploit this trait thoroughly. A good yo-yo puts as much weight as possible at the disc's \emph{rim}, not near its axle. A large moment of inertia lets it carry substantial angular momentum at modest speed, ``sleep'' at the end of its string, and sustain elaborate tricks. Whip-driven street tops exploit the other side: each lash applies a tangential force at the rim, far from the axis, creating a large torque that spins the top faster. Its rounded, balanced weight distribution keeps its mass even and its spin steady as a little stump. Then there is the long pole an acrobat spins on their head, with a basket or plate tied to its top. Placing mass high and far out deliberately increases the moment of inertia. The pole falls with a smaller angular acceleration, giving the performer time to step underneath and save it. Slow rotation can be a lifesaver.

\textbf{The fourth concept: angular momentum, rotation's quantity of motion.} In Chapter 10 we defined momentum as a quantity of motion and found it conserved in collisions. Rotation has a counterpart: \emph{angular momentum}, moment of inertia times angular velocity. Moment of inertia tells us how hard something is to set turning; angular velocity tells us how fast it turns now. Their product tells us how much rotation it carries. The deepest rotational law is this: \emph{when the external torques on a system sum to zero, its total angular momentum stays constant}. This is conservation of angular momentum, the rotational version of momentum conservation, equally reliable: no experiment to date has found an exception.

\textbf{A comparison table aligns the two ledgers.} Every rotational concept now has a relative in translation. This table is worth copying into your notebook. Almost everything left in this chapter, and almost every later rotation problem, amounts to translating between its two sides.
\end{mainline}

\begin{center}
\begin{tabular}{lll}
\toprule
\textbf{Role} & \textbf{Translation} & \textbf{Rotation} \\
\midrule
State: where now & Position & Angular position (radians) \\
State: rate of change & Velocity & Angular velocity \\
State: change of rate & Acceleration & Angular acceleration \\
Resistance to change & Mass & Moment of inertia \\
Ability to change motion & Force & Torque \\
``Quantity of motion'' & Momentum & Angular momentum \\
Law of motion change & \shortstack[l]{Net force changes\\momentum} & \shortstack[l]{External torque changes\\angular momentum} \\
Conservation & \shortstack[l]{No external force:\\momentum conserved} & \shortstack[l]{No external torque:\\angular momentum conserved} \\
\bottomrule
\end{tabular}
\end{center}

\step{The Model in One Sentence}

\begin{oneline}
Rotation changes in response to torque, force times lever arm, not force alone; moment of inertia measures resistance to turning and grows as mass moves farther from the axis; when external torques sum to zero, the system's angular momentum, moment of inertia times angular velocity, stays constant and is merely redistributed within it.
\end{oneline}

Remember that sentence and you have the key to eggs, skating, tops, and neutron stars alike.

\step{The Mathematical Translator}

\begin{mainline}
First translate the description. Angular position \(\theta\) in radians, angular velocity \(\omega\), and angular acceleration \(\alpha\) are defined just like their straight-line relatives, replacing distance with angle:
\[
  \omega=\frac{\Delta\theta}{\Delta t},\qquad
  \alpha=\frac{\Delta\omega}{\Delta t}.
\]
Read \(\omega\) as radians turned per second, and \(\alpha\) as the change in angular velocity per second. Routine checks: if \(\omega=0\), the object does not rotate. If \(\omega\) is large but \(\alpha=0\), it spins fast but \emph{uniformly}, like Earth. Negative \(\alpha\) need not mean slowing down: it depends on whether its sign matches that of \(\omega\). This is the same trap as mistaking negative acceleration for deceleration in straight-line motion.

Next, torque. For a lever arm \(d\), the perpendicular distance from the axis to the force's line of action, and a force \(F\):
\[
  \tau=Fd .
\]
Three routine checks. First, \(d=0\): a force passing through the axis produces no torque. Neither pushing right beside the hinge nor pushing along the door toward it will open the door, answering the second half of our third prediction question. Second, doubling the lever arm doubles \(\tau\). That explains longer wrenches for rusty nuts and why a small weight farther out on a balance can offset a large one. Third, direction: torques have signs, clockwise or counterclockwise. Two canceling torques leave rotation unchanged. Someone pushing from the other side with an opposing torque can hold the door in a stalemate.

Try real numbers. A bicycle pedal nut might require a tightening torque of about \(30\ \mathrm{N\cdot m}\). With a 25-centimeter wrench, you must apply \(30/0.25=120\) newtons, like lifting a 12-kilogram bucket of water. Not strong enough? Step on it: pressing part of your 60-kilogram body weight onto the wrench's end multiplies the force without changing the lever arm and readily supplies the torque. Conversely, a car manual may limit a delicate bolt to \(10\ \mathrm{N\cdot m}\). An experienced mechanic then \emph{holds the middle of the wrench}, deliberately shortening the arm to avoid accidentally snapping the bolt. The same formula amplifies effort one way and limits it the other.

Now translate moment of inertia. Imagine an object divided into small pieces, the \(i\)th having mass \(m_i\) at a distance \(r_i\) from the axis. Then
\[
  I=\sum_i m_i r_i^2 .
\]
This is the chapter's only formula with an inherent square, and deserves a close look. Squared distance makes moving mass outward the \emph{costliest} way to increase inertia. Check: if all the mass lies on the axis, \(r_i=0\), then \(I=0\); this ideal thin rod turns effortlessly about its axis. Moving the same mass twice as far out quadruples its contribution, exactly reproducing what your arms felt in the twisting experiment. Two useful results, derived in the mathematical bridge: for a \emph{thin ring} of mass \(M\) and radius \(R\) about its central axis, all the mass lies at distance \(R\), so \(I=MR^2\). A \emph{solid disc} of the same mass and radius has its mass much closer to the axis on average, giving \(I=\tfrac12 MR^2\). An equally large, equally heavy ring has twice the disc's moment of inertia. This difference will roll down a slope in our challenges.

\begin{center}
\begin{tikzpicture}[scale=0.85]
  % Thin ring: all mass at the rim
  \draw[very thick,red!70!black,fill=red!8] (0,0) circle (1.5);
  \draw[very thick,red!70!black,fill=white] (0,0) circle (1.25);
  \draw[-{Stealth[length=2.5mm]},thin] (0,0) -- (1.06,1.06) node[midway,above left=-2pt] {\(R\)};
  \node at (0,-1.95) {Thin ring: \(I=MR^{2}\)};
  % Solid disc: mass distributed throughout
  \begin{scope}[xshift=5.2cm]
    \draw[very thick,blue!70!black,fill=blue!10] (0,0) circle (1.5);
    \foreach \r in {0.45,0.9} \draw[thin,blue!40!black] (0,0) circle (\r);
    \draw[-{Stealth[length=2.5mm]},thin] (0,0) -- (1.06,1.06) node[midway,above left=-2pt] {\(R\)};
    \node at (0,-1.95) {Solid disc: \(I=\tfrac12 MR^{2}\)};
  \end{scope}
\end{tikzpicture}
\end{center}

Finally, our central quantity, angular momentum:
\[
  L=I\omega ,
\]
The conservation law says that \(L\) stays constant when external torques sum to zero. Comparing before and after gives \(I_1\omega_1=I_2\omega_2\). Check a limiting case: halve \(I_2\) and \(\omega_2\) doubles. Moment of inertia and angular velocity are opposite ends of a seesaw: when one goes down, the other must rise.
\end{mainline}

\textbf{An estimate with real numbers: figure skating.} With arms and one leg extended, a skater's moment of inertia is about \(3.2\ \mathrm{kg\cdot m^2}\): a body mass around 60 kilograms, mostly distributed twenty or thirty centimeters from the axis. Starting at one turn per second, or \(\omega_1=2\pi\ \mathrm{rad/s}\), she has angular momentum about \(3.2\times6.28\approx20\ \mathrm{kg\cdot m^2/s}\). Pulling her arms tightly against her body concentrates mass within a ten-centimeter radius, lowering her moment of inertia to about \(1.0\ \mathrm{kg\cdot m^2}\). Frictional torque from the ice is small, so angular momentum is approximately conserved over a few seconds:
\[
  \omega_2=\frac{I_1}{I_2}\omega_1\approx\frac{3.2}{1.0}\times 1\ \text{turn/s}\approx 3.2\ \text{turn/s}.
\]
That is just the order of magnitude of the sudden speedup seen in competition footage, the physics behind triple-and-a-half and quadruple jumps. One further account: pulling in raises rotational kinetic energy \(E=\tfrac12 I\omega^2\) from about \(63\) joules to about \(200\) joules. Conservation does not supply free energy. Her muscles do work pulling her arms inward against the outward-flinging tendency. Angular momentum and energy conservation keep separate accounts, just as in Chapters 9 and 10.

\begin{mathbridge}
Why radians rather than degrees? Only radians keep the arc-length formula clean: a point at distance \(r\) has linear speed \(v=\omega r\) and tangential acceleration \(a_{\text{tan}}=\alpha r\). Notice the implication: on one turntable, \(\omega\) is identical everywhere, while \(v\) grows in proportion to \(r\). Rim points run much faster than points near the center despite their identical angular velocities. This is the sharpest difference between rotation and translation: \(\omega\) is one number shared by the entire rigid body, whereas \(v\) differs from point to point.

The full torque formula is \(\tau=rF\sin\varphi\), where \(\varphi\) is the angle between the position vector and the force. It automatically includes both extremes: a force passing through the axis, with \(\varphi=0\) or \(\pi\), gives zero torque; a force perpendicular to the lever arm gives maximum torque. In vector form, \(\bm{\tau}=\bm r\times\bm F\), directed along the axis by the right-hand rule. Similarly, \(\bm L=I\bm\omega\) is a vector along the axis. Newton's second law for rotation is
\[
  \bm\tau_{\text{ext}}=\frac{d\bm L}{dt},
\]
Here the parallel between translation and rotation is complete: replace \(\bm F\) with \(\bm\tau\) and \(\bm p\) with \(\bm L\), and the form is identical. The disc's moment of inertia, \(I=\tfrac12MR^2\), also follows from one integral: slice the disc into infinitely many thin rings; a ring at radius \(r\) contributes \(r^2\,dm\). Integrate from \(0\) to \(R\).
\end{mathbridge}

\begin{challenge}
Why does a top circle and nod instead of falling? The simplified answer: gravity exerts a nonzero torque \(\bm\tau\) about the top's support point. By \(\bm\tau=d\bm L/dt\), it continually changes the direction, not the magnitude, of the angular-momentum \emph{vector}. The arrow \(\bm L\) is forced to sweep a cone uniformly around the vertical: precession. A rigorous treatment resolves angular velocity along three principal axes of inertia, using Euler's equations. It reveals that a freely rotating rigid body is stable about its largest or smallest inertia axis, but exhibits the famous flipping instability about its intermediate axis. Toss a spinning T-shaped wrench inside the International Space Station and it flips periodically. This requires vector calculus and is beyond our requirements, but remember the shape of the conclusion: \emph{torque changes where the angular-momentum arrow points}. That underlies the capabilities of bicycle wheels, gyroscopes, and gyrocompasses.
\end{challenge}

\step{The Model's Limits}

Every model must declare its territory. Four boundaries of the rotation model are especially easy to cross.

\textbf{Boundary 1: the rigid-body assumption.} Our formulas assume a \emph{rigid body}, an object of fixed shape whose parts rotate together. A boiled egg approximately qualifies; a raw egg does not. When its shell turns, sticky friction gradually brings the white and yolk along. Once internal parts can slip relative to one another, one \(\omega\) cannot describe the object: each part needs its own account. Cats righting themselves in the air and divers tucking and extending their bodies are advanced non-rigid-body maneuvers. Angular momentum is still conserved, but the calculations are much harder.

\textbf{Boundary 2: conservation requires zero external torque, not zero external force.} These are easily confused. An object can experience a sizable net force while conserving angular momentum: a flywheel on a frictionless fixed axle experiences a large support force, but its zero lever arm gives zero torque. Conversely, an object can spin faster with zero net force. Imagine two people on opposite sides of a door applying equal, opposite forces to the ends of the same-side handle? A clearer example is a steering wheel: one hand pushes up while the other pushes down. Net force is zero, torque is not, and the wheel turns. For rotation, always ask about torque first.

\textbf{Boundary 3: fixed and free axes.} We assumed a fixed axis. A free axis brings richer behavior: top precession, Earth's axis slowly circling over 26,000 years, and a tossed wrench flipping. The phenomena are striking, but the ledger's principle is unchanged. Only the geometry of which direction is conserved becomes more involved.

\textbf{Boundary 4: do not keep old misconceptions alive.} Two deserve special attention. First, the centrifugal force that supposedly throws you outward on a turn is not a new force. Tire--road friction supplies the centripetal force turning a car. A water droplet leaving a spinning umbrella tangentially is not flung by a new force; it loses the pull needed to keep moving in a circle. Like torque, centripetal force is a \emph{description of the role} existing forces, such as gravity, elasticity, and friction, play in a problem, not a new member of their family. Second, a widespread overexplanation deserves a counterexample: ``Bicycles stay upright entirely because their wheels are gyroscopes.'' The gyroscopic effect exists, but in 2011 researchers built a special bicycle with counter-rotating wheel pairs at front and rear, precisely canceling its total rotational angular momentum. It could still stabilize itself while moving. Bicycle stability actually comes from the front-axis geometry, with the tilted fork steering the front wheel toward a lean, together with the rider's continual adjustments. The lesson worth remembering most is this: \emph{a correct physical mechanism need not be the only mechanism at work}. Phone gyroscope chips and satellite attitude-control wheels, however, genuinely rely on angular momentum. A satellite turns without fuel using motor-driven flywheels inside it: as a wheel speeds up one way, the spacecraft slowly turns the other. Angular momentum transfers between them while the total remains unchanged. These are momentum wheels; the Hubble telescope uses them to aim at deep-space targets.

\step{Explaining the Opening Phenomenon}

Return to the two breakfast eggs and check each account with the new model.

Why does the raw egg spin poorly? Your fingers turn the \emph{shell}, giving it angular velocity. Its liquid white and yolk are dragged along only by viscous friction. The shell rotates faster than the contents: their angular velocities differ. This is not a rigid body, so Boundary 1 applies. Internal viscous torques continually transfer angular momentum from shell to liquid, while table friction drains angular momentum from the whole system. The egg soon slows. Much of the angular momentum you supplied dissipates into a tiny amount of heat inside the liquid.

Why does the boiled egg spin longer? Cooking solidifies white and yolk into a whole, approximately a rigid body. One flick gives inside and outside angular velocity together. There is no internal loss, only small resisting torques from the table and air. The egg holds its angular momentum tightly and naturally spins steadily for longer.

The best part is the revival. During the second you hold the boiled egg, friction from your finger stops its \emph{shell}. But the solidified white and yolk are not a perfect rigid body. The yolk, especially, retains a little room to slip against the white. At the instant you stop the shell, some internal rotation remains: not all the angular momentum has passed to your finger; some hides inside. After release, internal viscous friction works in reverse. The still-rotating contents drag the shell around again through internal frictional torque, restarting the whole egg. The account closes perfectly: its revived spin \emph{must be slower than before}, because your finger and the table removed some angular momentum while it was held. The first prediction answer is b. Careful observation reveals that a longer hold gives weaker revival. Hold it until the contents also stop completely and it cannot revive.

We have already done the skater's account. Pulling her arms toward the axis reduces moment of inertia to one third, and angular momentum conservation forces angular velocity to triple. The bicycle's answer is a combination of ``gyroscopic effect \(+\) fork geometry \(+\) rider adjustments.'' All three parts are indispensable; treating any one mechanism as the entire cause misreads the ledger.

Finally, why has Earth spun for over four billion years without stopping? The answer is now almost free. Earth is approximately a huge rigid body floating in vacuum. Solar and lunar gravity exert almost no net torque because their lines of action pass near its center. Its rotational angular momentum has nowhere to go, so it remains. Earth is not tireless; it needs no effort at all. Like uniform straight-line motion, rotation is a default state that needs neither explanation nor maintenance. For a finer account, lunar tides rubbing against the seabed act like an extremely gentle hand on Earth's boiled egg, removing a little rotational angular momentum each century. The day lengthens by about \(1.7\) milliseconds per century; over four billion years, it has grown from a few hours to twenty-four. Where that small outflow goes is the account pursued in Challenge 2.

\step{Thought Experiments and Challenges}

\begin{thought}[Can a Floating Astronaut Turn Around?]
An astronaut floats motionless and neutrally buoyant in the middle of a spacecraft, unable to touch any wall. Can she change her orientation \emph{by herself}, turning from facing the hatch to facing away?

Intuition says no: no pivot, no external force, and total angular momentum must remain zero. How could she turn? Yet she can. Here is the method cats use when falling: extend both arms left and twist the upper body right. To keep total angular momentum zero, the lower body and legs rotate left. Then draw in the arms, bend the legs into a different posture, and twist back. Because the upper and lower body's moments of inertia differ between extended and tucked positions, the two angular displacements do not cancel. Throughout the sequence, total angular momentum stays zero, yet her orientation changes. This is not magic but angular momentum conservation for a non-rigid body: conservation constrains the \emph{total}, not the \emph{pose}. Space-station astronauts really use similar movements to reorient themselves. In a world without external torque, the rule is not ``no turning'' but ``every bit of turning must have an opposite partner.''
\end{thought}

At the other end of the cosmic scale, angular momentum conservation writes a wilder story. A star of roughly the Sun's mass undergoes core collapse late in life. Its radius shrinks from about \(7\times10^5\) kilometers, the Sun's current radius, to about \(10\) kilometers, retaining almost all its mass. Moment of inertia scales with radius squared, so \(I\) falls to roughly \((10/7\times10^5)^2\approx2\times10^{-10}\) of its original value. If angular momentum is approximately conserved, its angular velocity must increase about five billionfold. The Sun turns once in about 25 days; the collapsed core, a \emph{neutron star}, could theoretically turn thousands of times per second. Observed millisecond pulsars spin more than 700 times per second at the fastest, the same order of magnitude as our estimate. Their equatorial surface speed approaches a fraction of the speed of light. They have simply taken the skater's arm-pulling trick to its extreme. From breakfast eggs to stellar corpses, the same law keeps the books.

\begin{puzzles}
\item Release a solid disc and a thin ring simultaneously at the top of the same slope. They have equal mass and radius, and both roll without slipping. Which reaches the bottom first? Hint: the gravitational potential energy is the same and becomes kinetic energy, but that energy divides into translation and rotation. A larger moment of inertia takes a larger rotational share, leaving less for translation. Recall \(I_{\text{ring}}=MR^2\) and \(I_{\text{disc}}=\tfrac12MR^2\). The answer depends on neither mass nor radius.
\item Astronomical measurements show that Earth's day lengthens by about \(1.7\) milliseconds per century while the Moon recedes about \(3.8\) centimeters per year. Why do these happen together? Hint: tidal friction is like a hand on Earth's boiled egg. The Earth--Moon system experiences approximately zero external torque. Where does Earth's lost spin angular momentum go? As the Moon moves farther away, does its orbital angular momentum increase or decrease?
\item After pulling in her arms to spin faster, a skater throws a glove \emph{tangentially}. Does her spin rate change greatly at that instant? Hint: the angular momentum the glove carries away depends on its motion. First decide whose angular momentum is conserved: hers alone, or hers plus the glove's? Then consider the difference if she simply \emph{lets go gently} instead of throwing it.
\end{puzzles}

We leave two threads for later. First, Chapter 58 will investigate the origin of angular momentum conservation. It is not an accidental empirical regularity but follows directly from a symmetry: all directions in space are equivalent. One of twentieth-century physics' deepest discoveries is that conservation laws all come from symmetries. Second, in the microscopic world angular momentum reveals a quantum face. Electrons and protons carry intrinsic angular momentum called spin, which cannot be understood as little balls literally rotating and comes in quantized amounts. That is Chapter 47's story. However deep we go, the ledger's format stays the same: first ask about the system's current rotational state, then what torque changes it, and finally how we measure it. Rotation and translation invite the same questions.
